----------------------------------------------------------------------
Report of the Referee -- DK13815/Heinzel
----------------------------------------------------------------------

This is a long, dense, and technically demanding paper on gravitational-wave
population inference. It will be accessible to a relatively small community of
specialists. That said, the content is extremely valuable and, in my view,
should be published. The paper carries out low-level and foundational
calculations that are currently scattered, implicit, or unavailable in the
literature, and presents a systematic analysis. Making these calculations
available in a refereed publication is a significant service to the field. I
expect this work to remain a reference for many years.

Here are some comments and suggestions.

- The title mentions only Bayesian statistics. While the paper is indeed heavily
statistical, its motivation, context, and examples are specific to
gravitational-wave astronomy. The title should reflect this more clearly. As it
stands, it risks misleading readers outside GW astronomy while underselling the
relevance to the GW community.
- The introduction is very LVK-centric; I encourage the authors to reframe it
within a broader GW context. For example, population inference for LISA is
likely to look very different in the context of a global fit, where multiple
sources overlap and sources from the same population may be either individually
resolved or contribute to a stochastic foreground. Mentioning such scenarios
would help position the paper as a general contribution rather than an LVK
technical fix.
- Similarly, I found that the introduction (perhaps implicitly) suggests that
the issues tackled here are intrinsic shortcomings of hierarchical Bayesian
statistics, which is misleading. Rather, these issues arise from how
hierarchical Bayesian inference is implemented in LVK analyses. In particular,
they stem from the now-standard LVK separation between searches, parameter
estimation, and population inference. This separation forces population analyses
to recycle posterior samples from parameter estimation and to rely on injection
campaigns to account for selection effects. The resulting statistical
pathologies are therefore implementation-driven, not fundamental. I encourage
the author to clarify this point.
- Page 2 surveys previous LVK-related work. I would like to point out the
related discussion in arXiv:2502.12156, where the authors similarly conclude
that conditioning on the variance of the likelihood is more stringent than using
the effective number of samples.
- Throughout the paper, there is an almost obsessive referencing to Farr,
Essick, Talbot, and Golomb. Those papers are certainly relevant and should be
cited, but the present study provides a substantially new contribution and
should stand on its own. The authors do not need to be quite so deferential to
earlier work.
- I have a conceptual question regarding the proposed correction to the
likelihood, specifically when going from Eq. 18 to Eq. 19. The correction
assumes log-normality of the estimator, which is only justified when the number
of samples M is large. However, while the estimator is biased at finite M, it is
asymptotically unbiased, with the bias vanishing as M goes to infinity.
Therefore, in the regime where the log-normal approximation is valid, the bias
is already negligible. What, then, is the practical role of this correction? Is
this circular reasoning?
- I find footnote 6 unclear. Are the authors implying that there is a mistake in
Farr [56]? If so, this should be stated clearly. Leaving the issue ambiguous is
likely to confuse future readers of both papers.
- Has the likelihood correction proposed here (Eq. 24) been implemented in any
LVK population studies to date?
- I have some confusion (extending beyond this paper, actually) regarding the
common practice of thresholding on uncertainty in GW population inference,
discussed here in Sec. 5.A. These practices, which are widely adopted, are based
either on the variance of the likelihood or on the effective number of samples.
When implemented, they modify the likelihood (some GW literature has referred to
this as a data-dependent prior, which I believe is incorrect). Given this, how
can one be confident that the *thresholded* likelihood is unbiased? The
likelihood and posterior bias calculations presented in this paper are performed
on the *unthresholded* likelihood estimator. Do these calculations remain valid
once thresholding is applied?
- Starting from the abstract, the paper appears to recommend adopting a generic
condition  $\hat{E} < 0.2$ . However, in the final paragraph of Sec. 6 the
authors correctly note that the dependence on the population model is complex,
and that what is valid for one model may not be valid for another. I agree with
this latter statement. Given that, I do not understand why a universal numerical
threshold of 0.2 is recommended. This should be clarified.
- Finally, I personally found Appendix B more confusing than helpful.

Overall, this is an important and technically solid paper that deserves
publication. Most of my comments concern framing and clarity rather than
correctness. Addressing them would improve the paper's accessibility and
long-term impact.